\chapter{Forma bilineal electrogravitatoria y ecuaciones de respuesta cruzada}
\label{chap:forma-electrogravitatoria}

\section{Coordenadas internas de fuente}

Sea \(Z\) la secci\'on de impedancia del vac\'io y sea \((Q,\ell)\) la pareja
del cap\'itulo anterior. A una carga \(q\) y una energ\'ia \(E\) asociamos

\[
\xi_q(q)=\sqrt{\frac{Z}{4\pi\hbar}}\,q,
\qquad
\xi_E(E)=\frac{E}{Q}.
\]

El vector interno de fuente es

\[
X(q,E)=
\begin{pmatrix}\xi_q(q)\\ \xi_E(E)\end{pmatrix},
\qquad
\eta_{1,1}=
\begin{pmatrix}1&0\\0&-1\end{pmatrix}.
\]

La forma \(X_1^{\mathsf T}\eta_{1,1}X_2\) separa orientaci\'on el\'ectrica y
cierre gravitatorio dentro de un mismo plano adimensional.

\section{Unificaci\'on est\'atica de los potenciales}

\begin{theorem}[Forma electrogravitatoria]
\label{thm:rm-emg-potential}
Si \(Zc=\varepsilon^{-1}\), \(E_i=m_ic^2\) y
\(Q\ell=\hbar c\), entonces
\begin{equation}
\boxed{
V_{12}(r)=\frac{Q\ell}{r}
X(q_1,E_1)^{\mathsf T}\eta_{1,1}X(q_2,E_2)}
\label{eq:rm-emg-potential}
\end{equation}
es exactamente la suma del potencial de Coulomb y el potencial de Newton.
\end{theorem}

\begin{proof}
El t\'ermino el\'ectrico es
\[
\frac{Q\ell}{r}\frac{Zq_1q_2}{4\pi\hbar}
=\frac{Zcq_1q_2}{4\pi r}
=\frac{q_1q_2}{4\pi\varepsilon r}.
\]
El t\'ermino gravitatorio es
\[
-\frac{Q\ell}{r}\frac{E_1E_2}{Q^2}
=-\frac{\ell}{Q}\frac{E_1E_2}{r}
=-\frac{G}{c^4}\frac{m_1m_2c^4}{r}
=-\frac{Gm_1m_2}{r}.
\]
\end{proof}

La fuerza radial asociada adopta la forma

\begin{equation}
F_{12}(r)=\frac{Q}{\ell}
\left(\frac{\ell}{r}\right)^2
X_1^{\mathsf T}\eta_{1,1}X_2.
\label{eq:rm-emg-force}
\end{equation}

El factor \(Q/\ell=\mathscr T_P\) es la tensi\'on universal; la forma
adimensional decide la intensidad y el signo de su publicaci\'on.

\section{Cono de extremalidad de las respuestas constitutivas y clasificación de sus rayos}

Para una fuente individual, el discriminante

\[
\Delta_{\rm EMG}(q,E)=E^2-\mathcal Q_e(q)^2,
\qquad
\mathcal Q_e(q)=Q\sqrt{\frac{Z}{4\pi\hbar}}\,q,
\]

induce los tres sectores

\[
\Delta_{\rm EMG}>0,
\qquad
\Delta_{\rm EMG}=0,
\qquad
\Delta_{\rm EMG}<0.
\]

La condici\'on nula equivale a

\begin{equation}
E^2=\mathcal Q_e(q)^2
\quad\Longleftrightarrow\quad
Gm^2=\frac{q^2}{4\pi\varepsilon}.
\label{eq:rm-rn-extremality}
\end{equation}

Es la condici\'on de extremalidad electrogravitatoria en la normalizaci\'on de
Reissner--Nordstr\"om. El trit adquiere aqu\'i una realizaci\'on est\'atica:
subextremalidad, umbral extremo y superextremalidad son los tres signos de una
misma forma cuadr\'atica.

Para la carga elemental construida por

\[
e^2=\frac{4\pi\alpha_{\rm HMT}\hbar}{Z},
\]

la coordenada satisface

\begin{equation}
\boxed{\xi_q(e)^2=\alpha_{\rm HMT}.}
\label{eq:rm-electron-electric-coordinate}
\end{equation}

Esta igualdad sit\'ua la carga central en el plano de fuentes sin utilizar la
masa electr\'onica ni importar su teor\'ia completa.

\section{Coeficiente de Einstein--Maxwell}

Escribamos

\[
\kappa=\frac{8\pi G}{c^4}=8\pi\mathscr C_P.
\]

En una convenci\'on donde el tensor electromagn\'etico se expresa como
\(T_{\mu\nu}^{\rm EM}=\mu^{-1}\mathcal Q_{\mu\nu}(F,g)\), el coeficiente
gravitatorio de la forma cuadr\'atica es \(\kappa/\mu\). Usando

\[
\ell^2=\frac{G\hbar}{c^3},
\qquad
Z=\mu c,
\qquad
e^2=\frac{4\pi\alpha\hbar}{Z},
\]

se obtiene

\begin{equation}
\boxed{
\frac{\hbar^2}{8\pi\ell^2}
\frac{(\kappa/\mu)}{e^2}
=\frac{1}{4\pi\alpha_{\rm HMT}}.}
\label{eq:rm-einstein-maxwell-alpha}
\end{equation}

\begin{proof}
Primero,
\[
\frac{\kappa}{\mu}
=\frac{8\pi G}{c^4\mu}
=\frac{8\pi\ell^2}{\hbar\mu c}
=\frac{8\pi\ell^2}{\hbar Z}.
\]
Al dividir por \(e^2=4\pi\alpha\hbar/Z\) y multiplicar por
\(\hbar^2/(8\pi\ell^2)\) queda \(1/(4\pi\alpha)\).
\end{proof}

En la interfaz electrod\'ebil racionalizada,

\[
g^2=\frac{4\pi\alpha}{1-D_A},
\qquad
g'^2=\frac{4\pi\alpha}{D_A},
\]

y por ello

\begin{equation}
\boxed{
\frac{\hbar^2}{8\pi\ell^2}
\frac{(\kappa/\mu)}{e^2}
=\frac1{g^2}+\frac1{g'^2}.}
\label{eq:rm-cross-einstein-ew}
\end{equation}

La igualdad obliga a conmutar los lectores de gravedad, vac\'io, carga y
mezcla gauge. Dentro del grafo actual, \(e,g,g'\) descienden de
\(\alpha_{\rm HMT}\); la ecuaci\'on es un control cruzado y no una segunda
derivaci\'on independiente de \(\alpha\).

\section{Polígono de publicaciones de la estructura fina}

El macroobjeto dodecaf\'asico se publica mediante el sistema

\begin{equation}
\boxed{
\alpha_{\rm HMT}
=\frac{Ze^2}{4\pi\hbar}
=\frac{ZG_0}{4}
=\frac{Z}{2R_K}
=-\frac{9}{100\pi}\log D_A
=\frac{(1-D_A)g^2}{4\pi}
=\frac{D_Ag'^2}{4\pi}.}
\label{eq:rm-alpha-polygon}
\end{equation}

La primera construcci\'on causal de \(\alpha\) sigue siendo la rueda firmada,
el sello y el acarreo. Las igualdades posteriores son publicaciones
metrol\'ogicas o controles de compatibilidad. La eliminaci\'on de variables
produce

\[
(1-D_A)g^2=D_Ag'^2
=\frac{Ze^2}{\hbar}
=\pi G_0Z
=\frac{2\pi Z}{R_K}.
\]

La forma dual es

\[
\frac1{g^2}+\frac1{g'^2}
=\frac{\hbar}{Ze^2}
=\frac{R_K}{2\pi Z}
=\frac1{\pi G_0Z}.
\]

Cada ciclo metrol\'ogico debe marcarse como \emph{round trip}: si \(e\),
\(R_K\) o \(G_0\) fueron definidos desde \(\alpha\), su sustituci\'on inversa
comprueba coherencia, pero no a\~nade informaci\'on selectiva.

\section{Covariancia de la inversi\'on de acci\'on}

Bajo la transformaci\'on bidireccional,

\[
\frac{e_{\rm pre}^2}{e_{\rm ret}^2}=R_{\rm act},
\qquad
\frac{G_{\rm pre}}{G_{\rm ret}}=R_{\rm act}^{-1}.
\]

Por tanto,

\begin{equation}
\boxed{G_{\sigma}e_{\sigma}^2=\text{constante de hoja}.}
\label{eq:rm-ge2-invariant}
\end{equation}

En forma constitutiva,

\[
\frac{Ze_{\sigma}^2G_{\sigma}}
{c^5t_0^2}
=4\pi\theta_{108}\alpha_{\rm HMT},
\qquad
\theta_{108}=\left(\frac{54}{\pi}\right)^2.
\]

La carga conserva la orientaci\'on de la acci\'on y la respuesta gravitatoria
la compensa; su producto pertenece a la geometr\'ia compartida.

\Needspace{12\baselineskip}
\section{Representación tensorial del portador gravitatorio y reducción transversal de sus grados de libertad}

La forma bilineal \eqref{eq:rm-emg-potential} cierra el sector est\'atico. La
promoci\'on a una din\'amica completa exige un realizador

\[
\mathcal R_{\rm EM}:
\mathcal X_{\TPK}^{\rm enr}
\longrightarrow
(e^I{}_{\mu},\omega^{IJ}{}_{\mu},A_\mu)
\]

que preserve la identidad de Bianchi, las ecuaciones de Maxwell, la
conservaci\'on de corriente y los t\'erminos de borde. La identidad
\eqref{eq:rm-cross-einstein-ew} fija el coeficiente que ese realizador debe
publicar; no reemplaza sus ecuaciones diferenciales.

\begin{falsifierbox}
Refutan la forma electrogravitatoria: un coeficiente distinto en Coulomb o
Newton al expandir \eqref{eq:rm-emg-potential}; un cono nulo que no reproduzca
\eqref{eq:rm-rn-extremality}; la p\'erdida de
\eqref{eq:rm-einstein-maxwell-alpha}; o la utilizaci\'on de una publicaci\'on
descendiente de \(\alpha\) como selector retroactivo de su rama dodecaf\'asica.
\end{falsifierbox}

\paragraph{Estructura causal de la confluencia electrogravitatoria.}
La composición demostrada en este capítulo se organiza en cinco etapas
tipadas. El estado enriquecido suministra primero la fase
\(\alpha_{\rm HMT}\) y su orientación. El operador constitutivo publica después
\(e\), \(g\) y \(g'\) en sus codominios respectivos. El cierre gravitatorio
determina \(G\) y la cronología \(\theta_{108}\). Las identidades
\eqref{eq:rm-emg-potential}--\eqref{eq:rm-cross-einstein-ew} coordinan esas
publicaciones mediante un coeficiente común. Finalmente,
\(\mathcal R_{\rm EM}\) transporta el resultado hacia los campos
\((e^I{}_{\mu},\omega^{IJ}{}_{\mu},A_\mu)\), preservando las identidades
diferenciales y los términos de frontera. Esta factorización fija la
antecedencia de la construcción HMT respecto de su realización tensorial y
prepara las relaciones constitutivas desarrolladas en el capítulo siguiente.
